% Values of the integral \int M_u M dl (after Ghali & Neville 1997)
% Small shape macros (M diagrams over length l)
\newcommand{\shpRect}[1]{\begin{tikzpicture}[scale=0.55,baseline=(current bounding box.center),font=\tiny]
  \draw[thick] (0,0) rectangle (2.4,0.7); \node[right] at (2.4,0.35) {$#1$}; \draw[dim] (0,-0.35) -- (2.4,-0.35) node[midway,fill=white,inner sep=0.5pt] {$l$};\end{tikzpicture}}
\newcommand{\shpTriR}[1]{\begin{tikzpicture}[scale=0.55,baseline=(current bounding box.center),font=\tiny]
  \draw[thick] (0,0) -- (2.4,0) -- (2.4,0.8) -- cycle; \node[right] at (2.4,0.4) {$#1$}; \draw[dim] (0,-0.35) -- (2.4,-0.35) node[midway,fill=white,inner sep=0.5pt] {$l$};\end{tikzpicture}}
\newcommand{\shpTriL}[1]{\begin{tikzpicture}[scale=0.55,baseline=(current bounding box.center),font=\tiny]
  \draw[thick] (0,0) -- (2.4,0) -- (0,0.8) -- cycle; \node[left] at (0,0.4) {$#1$}; \draw[dim] (0,-0.35) -- (2.4,-0.35) node[midway,fill=white,inner sep=0.5pt] {$l$};\end{tikzpicture}}
\newcommand{\shpTrap}[2]{\begin{tikzpicture}[scale=0.55,baseline=(current bounding box.center),font=\tiny]
  \draw[thick] (0,0) -- (2.4,0) -- (2.4,1.0) -- (0,0.45) -- cycle; \node[left] at (0,0.25) {$#1$}; \node[right] at (2.4,0.6) {$#2$}; \draw[dim] (0,-0.35) -- (2.4,-0.35) node[midway,fill=white,inner sep=0.5pt] {$l$};\end{tikzpicture}}
\newcommand{\shpPeak}[1]{\begin{tikzpicture}[scale=0.55,baseline=(current bounding box.center),font=\tiny]
  \draw[thick] (0,0) -- (2.4,0) -- (0.9,0.9) -- cycle; \draw (0.9,0) -- (0.9,0.9); \node[right] at (0.9,0.45) {$#1$};
  \draw[dim] (0,1.2) -- (0.9,1.2) node[midway,fill=white,inner sep=0.5pt] {$\alpha l$}; \draw[dim] (0.9,1.2) -- (2.4,1.2) node[midway,fill=white,inner sep=0.5pt] {$\beta l$};
  \draw[dim] (0,-0.35) -- (2.4,-0.35) node[midway,fill=white,inner sep=0.5pt] {$l$};\end{tikzpicture}}
\newcommand{\shpPara}[1]{\begin{tikzpicture}[scale=0.55,baseline=(current bounding box.center),font=\tiny]
  \draw[thick] (0,0) -- (2.4,0); \draw[thick,domain=0:2.4,samples=30] plot (\x,{0.9*(1-((\x-1.2)/1.2)^2)}); \draw (1.2,0) -- (1.2,0.9); \node[right] at (1.2,0.45) {$#1$}; \node at (2.6,0.9) {$*$};
  \draw[dim] (0,-0.35) -- (2.4,-0.35) node[midway,fill=white,inner sep=0.5pt] {$l$};\end{tikzpicture}}
\newcommand{\shpParaA}[1]{\begin{tikzpicture}[scale=0.55,baseline=(current bounding box.center),font=\tiny]
  \draw[thick] (0,0) -- (2.4,0) -- (2.4,0.9); \draw[thick,domain=0:2.4,samples=30] plot (\x,{0.9*(\x/2.4)^2}); \draw[thin] (-0.6,0) -- (0.8,0); \node[right] at (2.4,0.45) {$#1$}; \node at (2.6,1.0) {$*$};
  \node[above,font=\tiny] at (0.1,0.05) {tangent}; \draw[dim] (0,-0.35) -- (2.4,-0.35) node[midway,fill=white,inner sep=0.5pt] {$l$};\end{tikzpicture}}
\newcommand{\shpParaB}[1]{\begin{tikzpicture}[scale=0.55,baseline=(current bounding box.center),font=\tiny]
  \draw[thick] (0,0) -- (2.4,0) -- (2.4,0.9); \draw[thick,domain=0:2.4,samples=30] plot (\x,{0.9*(1-(1-\x/2.4)^2)}); \draw[thin] (1.4,0.9) -- (3.0,0.9); \node[right] at (2.4,0.45) {$#1$}; \node at (2.6,1.15) {$*$};
  \node[above,font=\tiny] at (0.5,0.55) {tangent}; \draw[dim] (0,-0.35) -- (2.4,-0.35) node[midway,fill=white,inner sep=0.5pt] {$l$};\end{tikzpicture}}
{\footnotesize
\renewcommand{\arraystretch}{1.6}
\begin{tabular}{|c|c|c|c|c|c|}\hline
  \diagbox[width=2.6cm]{$M$}{$M_u$} & \shpRect{b} & \shpTriR{b} & \shpTriL{b} & \shpTrap{b_1}{b_2} & \shpPeak{b}\\\hline
  \shpRect{a} & $abl$ & $\frac12abl$ & $\frac12abl$ & $\frac{al}{2}(b_1+b_2)$ & $\frac12abl$\\\hline
  \shpTrap{a_1}{a_2} & $\frac{bl}{2}(a_1+a_2)$ & $\frac{bl}{6}(a_1+2a_2)$ & $\frac{bl}{6}(2a_1+a_2)$ & \makecell{$\frac l6(2a_1b_1+a_1b_2+$\\$a_2b_1+2a_2b_2)$} & \makecell{$\frac{bl}{6}[(1+\beta)a_1$\\$+(1+\alpha)a_2]$}\\\hline
  \shpPeak{a} & $\frac12abl$ & $\frac{abl}{6}(1+\alpha)$ & $\frac{abl}{6}(1+\beta)$ & \makecell{$\frac{al}{6}[(1+\beta)b_1+$\\$(1+\alpha)b_2]$} & $\frac13abl$\\\hline
  \shpPara{a} & $\frac23abl$ & $\frac13abl$ & $\frac13abl$ & $\frac13al(b_1+b_2)$ & $\frac{abl}{3}(1+\alpha\beta)$\\\hline
  \shpParaA{a} & $\frac13abl$ & $\frac14abl$ & $\frac1{12}abl$ & $\frac{al}{12}(b_1+3b_2)$ & $\frac{abl}{12}(1+\alpha+\alpha^2)$\\\hline
  \shpParaB{a} & $\frac23abl$ & $\frac5{12}abl$ & $\frac14abl$ & $\frac{al}{12}(3b_1+5b_2)$ & $\frac{abl}{12}(5-\beta-\beta^2)$\\\hline
\end{tabular}\\[2pt]
{\scriptsize $*$ Second-degree parabola.}}
